\chapter{Pecahan: Keempat Operasinya}\label{ch:g8:fractions}

Bab ini melengkapi aritmetika \omterm{def:g6:fractions:def}{pecahan}: penjumlahan dan pengurangan
dengan \omterm{def:g4:fractions:def}{penyebut} \emph{apa pun}, perkalian, dan --- pendatang barunya
--- pembagian, yang ternyata perkalian dalam samaran. Tanda dari
\cref{ch:g8:negprod} ikut serta di dalamnya.

\section{Menjumlahkan dengan penyebut apa pun}

\begin{method}[Penyebut yang sama]\label{met:g8:fractions:add}
Untuk menjumlahkan atau mengurangkan $\frac ab$ dan $\frac cd$:
\begin{enumerate}
  \item carilah \omterm{def:g5:division:multiple}{kelipatan} bersama dari $b$ dan $d$ (\omterm{def:g2:mult:def}{hasil kali}
        $b \times d$ selalu berhasil; yang lebih kecil menghemat
        tenaga);
  \item tulislah ulang kedua \omterm{def:g6:fractions:def}{pecahannya} dengan \omterm{def:g4:fractions:def}{penyebut} itu;
  \item jumlahkan atau kurangkan \omterm{def:g4:fractions:def}{pembilangnya}; lalu sederhanakan.
\end{enumerate}
\end{method}

\begin{example}\label{ex:g8:fractions:add}
Hitunglah $\dfrac{5}{6} + \dfrac{3}{4}$. \omterm{def:g5:division:multiple}{Kelipatan} bersama dari $6$
dan $4$ adalah $12$:
\[
\frac{5}{6} = \frac{10}{12},
\qquad
\frac{3}{4} = \frac{9}{12},
\qquad
\frac{5}{6} + \frac{3}{4} = \frac{10 + 9}{12} = \frac{19}{12}.
\]
Dengan tanda: $\dfrac{2}{3} - \dfrac{7}{5} = \dfrac{10}{15} -
\dfrac{21}{15} = -\dfrac{11}{15}$.
\end{example}

\section{Mengalikan pecahan}

\begin{theorem}[Hasil kali pecahan]\label{thm:g8:fractions:mult}
\[
\frac{a}{b} \times \frac{c}{d} = \frac{a \times c}{b \times d} .
\]
\end{theorem}

\begin{proof}[Mengapa, lewat sebuah gambar]
Ambillah $\frac34$ dari $\frac25$ sebuah persegi: potonglah perseginya
menjadi $5$ jalur tegak lalu simpan $2$; potonglah mendatar menjadi
$4$ lalu simpan $3$ baris dari yang tersisa. Bagian yang disimpan
adalah kisi $3 \times 2$ sel kecil dari $4 \times 5$: \omterm{def:g6:fractions:def}{pecahannya}
$\frac{6}{20}$.
\end{proof}

\begin{omfigure}
\begin{tikzpicture}[scale=0.68]
  \draw[gray!50, thin] (0,0) grid[xstep=1, ystep=1] (5,4);
  \fill[omDef!18] (0,0) rectangle (2,4);
  \fill[omThm!35] (0,0) rectangle (2,3);
  \draw[thick] (0,0) rectangle (5,4);
  \draw[omDef, thick] (2,0) -- (2,4);
  \draw[omThm, thick] (0,3) -- (5,3);
  \node[below, font=\small, omDef] at (1,0) {$\frac25$ disimpan};
  \node[right, font=\small, omThm] at (5.05,1.5) {$\frac34$ darinya};
\end{tikzpicture}

{\small $\frac34 \times \frac25$: daerah yang terarsir ganda adalah
$3 \times 2 = 6$ sel dari $4 \times 5 = 20$, yaitu
$\frac{6}{20} = \frac{3}{10}$ --- \omterm{def:g4:fractions:def}{pembilangnya} dikalikan, \omterm{def:g4:fractions:def}{penyebutnya}
dikalikan.}
\end{omfigure}

\begin{example}\label{ex:g8:fractions:mult}
Sederhanakan \emph{sebelum} mengalikan, dengan mencoret faktor bersamanya:
\[
\frac{7}{12} \times \frac{9}{14}
= \frac{7 \times 9}{12 \times 14}
= \frac{7 \times 3 \times 3}{3 \times 4 \times 2 \times 7}
= \frac{3}{8}
\]
(faktor $7$ dan satu faktor $3$ saling meniadakan antara atas dan
bawah). Dengan tanda, aturan \cref{thm:g8:negprod:rules} berlaku lebih
dahulu: $\left(-\frac23\right) \times \left(-\frac{5}{4}\right) =
+\frac{10}{12} = \frac56$.
\end{example}

\section{Membagi pecahan}

\begin{definition}[Kebalikan]\label{def:g8:fractions:inverse}
\emph{Kebalikan}\index{kebalikan} sebuah bilangan $x$ yang bukan \omterm{def:g1:counting:zero}{nol}
adalah bilangan yang kalau dikalikan dengan $x$ memberi $1$.
Kebalikan $\frac ab$ adalah $\frac ba$, karena
$\frac ab \times \frac ba = \frac{ab}{ab} = 1$; kebalikan bilangan
bulat $n$ adalah $\frac1n$.
\end{definition}

\begin{theorem}[Membagi berarti mengalikan dengan kebalikannya]\label{thm:g8:fractions:div}
Untuk $c \neq 0$:
\[
\frac{a}{b} \div \frac{c}{d} = \frac{a}{b} \times \frac{d}{c} .
\]
\end{theorem}

\begin{proof}
Menurut definisinya, \omterm{def:g3:division:remainder}{hasil bagi} $q$ dari $\frac ab$ oleh $\frac cd$
adalah bilangan dengan $q \times \frac cd = \frac ab$. Periksalah
calonnya $q = \frac ab \times \frac dc$:
\[
\left(\frac ab \times \frac dc\right) \times \frac cd
= \frac ab \times \left(\frac dc \times \frac cd\right)
= \frac ab \times 1 = \frac ab . \qedhere
\]
\end{proof}

\begin{example}\label{ex:g8:fractions:div}
\[
\frac{3}{5} \div \frac{7}{10}
= \frac{3}{5} \times \frac{10}{7}
= \frac{3 \times 10}{5 \times 7}
= \frac{30}{35} = \frac{6}{7}.
\]
Pemeriksaan yang masuk akal: membagi dengan $\frac{7}{10}$, bilangan
yang lebih kecil daripada $1$, harus memberi hasil yang \emph{lebih
besar} daripada $\frac35$ --- dan memang $\frac67 > \frac35$
($\frac{30}{35} > \frac{21}{35}$).
\end{example}

\begin{example}[Garis pecahan di dalam garis pecahan]\label{ex:g8:fractions:complex}
Sebuah ``\omterm{def:g6:fractions:def}{pecahan} bertingkat'' hanyalah pembagian yang ditulis tegak:
\[
\frac{\ \frac{2}{3}\ }{\ \frac{5}{6}\ }
= \frac{2}{3} \div \frac{5}{6}
= \frac{2}{3} \times \frac{6}{5}
= \frac{12}{15} = \frac{4}{5}.
\]
Carilah garis \emph{utamanya} lebih dahulu (yang paling panjang), lalu
terapkan \cref{thm:g8:fractions:div}.
\end{example}

\begin{method}[Perhitungan campuran]\label{met:g8:fractions:mixed}
Pada bentuk yang mencampur \omterm{def:g6:fractions:def}{pecahan} dan keempat operasinya:
\begin{enumerate}
  \item terapkan prioritas yang biasa (\cref{ch:g7:priorities});
  \item ubahlah setiap pembagian menjadi perkalian dengan
        \omterm{def:g8:fractions:inverse}{kebalikannya};
  \item tentukan tandanya lebih dahulu
        (\cref{met:g8:negprod:compute}), lalu kalikan atau carilah
        \omterm{def:g4:fractions:def}{penyebut} yang sama;
  \item sederhanakan pada kesempatan yang paling awal.
\end{enumerate}
\end{method}

\begin{example}\label{ex:g8:fractions:mixed}
\begin{align*}
\frac12 + \frac{2}{3} \times \frac{9}{4}
&= \frac12 + \frac{2 \times 9}{3 \times 4}
&& \text{(hasil kalinya dahulu!)} \\
&= \frac12 + \frac{18}{12} = \frac12 + \frac32
= \frac{4}{2} = 2 .
\end{align*}
\end{example}

\section{Latihan}

\begin{exercise}[$\star$]\label{exo:g8:fractions:1}
Hitunglah lalu sederhanakan:
\[
\frac{3}{4} + \frac{2}{5}, \qquad
\frac{5}{6} - \frac{3}{8}, \qquad
\frac{7}{10} + \frac{8}{15} .
\]
\end{exercise}

\begin{exercise}[$\star$]\label{exo:g8:fractions:2}
Hitunglah (tandanya lebih dahulu):
\[
-\frac{2}{7} + \frac{3}{14}, \qquad
\frac{1}{6} - \frac{5}{4}, \qquad
-\frac{3}{5} - \frac{1}{2} .
\]
\end{exercise}

\begin{exercise}[$\star$]\label{exo:g8:fractions:3}
Hitunglah, dengan menyederhanakan sebelum mengalikan:
\[
\frac{5}{8} \times \frac{4}{15}, \qquad
\frac{9}{14} \times \frac{7}{3}, \qquad
\left(-\frac{6}{5}\right) \times \frac{10}{9} .
\]
\end{exercise}

\begin{exercise}[$\star$]\label{exo:g8:fractions:4}
Sebutkan \omterm{def:g8:fractions:inverse}{kebalikan} dari: $\dfrac{3}{7}$; \quad $5$; \quad
$\dfrac{1}{4}$; \quad $-\dfrac{2}{9}$.
\end{exercise}

\begin{exercise}[$\star$]\label{exo:g8:fractions:5}
Hitunglah:
\[
\frac{4}{9} \div \frac{2}{3}, \qquad
\frac{7}{5} \div 14, \qquad
6 \div \frac{3}{4} .
\]
\end{exercise}

\begin{exercise}[$\star$]\label{exo:g8:fractions:6}
Hitunglah \omterm{def:g6:fractions:def}{pecahan} bertingkatnya:
\[
\frac{\ \frac{3}{4}\ }{\ \frac{9}{8}\ },
\qquad
\frac{\ \frac{5}{6}\ }{\ 10\ },
\qquad
\frac{\ 2\ }{\ \frac{4}{7}\ } .
\]
\end{exercise}

\begin{exercise}[$\star$]\label{exo:g8:fractions:7}
Hitunglah, dengan mematuhi prioritasnya:
\[
\frac{1}{3} + \frac{1}{2} \times \frac{4}{5},
\qquad
\left(\frac{1}{3} + \frac{1}{2}\right) \times \frac{4}{5} .
\]
\end{exercise}

\begin{exercise}[$\star\star$]\label{exo:g8:fractions:8}
Sebuah tangki terisi $\frac{2}{5}$. Lalu ditambahkan $\frac{1}{3}$
\emph{dari \omterm{def:g2:measure:units}{daya tampung} tangkinya}. Berapa bagian tangkinya yang kini
terisi? Berapa bagian yang masih kosong?
\end{exercise}

\begin{exercise}[$\star\star$]\label{exo:g8:fractions:9}
Tiga perempat murid sebuah kelas lulus sebuah ujian; di antara mereka,
dua pertiganya memperoleh lebih dari $14$ dari $20$. Berapa bagian
kelasnya yang memperoleh lebih dari $14$? Kalau itu mewakili $12$
murid, berapa besar kelasnya?
\end{exercise}

\begin{exercise}[$\star\star$]\label{exo:g8:fractions:10}
Tali sepanjang $\frac{15}{2}$ m harus dipotong menjadi potongan
sepanjang $\frac{3}{4}$ m masing-masing. Berapa potongan yang
diperoleh? (Pembagian \omterm{def:g6:fractions:def}{pecahan} --- periksalah bahwa jawabannya
bilangan bulat.)
\end{exercise}

\begin{exercise}[$\star\star$]\label{exo:g8:fractions:11}
Hitunglah
\[
E = \frac{2}{3} - \frac{5}{6} \div \frac{10}{9},
\qquad
F = \left(-\frac{1}{2}\right)^2 \times \frac{8}{3} .
\]
\end{exercise}

\begin{exercise}[$\star\star\star$]\label{exo:g8:fractions:12}
Sederhanakanlah bentuk
\[
1 + \cfrac{1}{1 + \cfrac{1}{1 + 1}}
\]
(mulailah dari \omterm{def:g6:fractions:def}{pecahan} yang paling dalam lalu bekerjalah ke luar,
satu garis pada satu waktu).
\end{exercise}

\section{Soal: Pecahan Mesir}

\begin{problem}[{Soal akhir pekan --- setiap pecahan adalah jumlah
pecahan satuan yang berbeda-beda}]
\label{pb:g8:fractions:1}
\emph{\omterm{def:g6:fractions:def}{Pecahan} satuan} adalah \omterm{def:g6:fractions:def}{pecahan} yang \omterm{def:g4:fractions:def}{pembilangnya} $1$. Juru tulis
Mesir kuno --- papirus Rhind disalin sekitar tahun
$1550$~SM --- menulis setiap \omterm{def:g6:fractions:def}{pecahan} sebagai \omterm{def:g1:addition:def}{jumlah} \omterm{def:g6:fractions:def}{pecahan} satuan
yang \emph{berbeda-beda}: tidak pernah $\frac27 = \frac17 + \frac17$,
melainkan $\frac27 = \frac14 + \frac{1}{28}$. \omterm{def:g1:addition:def}{Jumlah} seperti itu
disebut \emph{tulisan Mesir} bagi \omterm{def:g6:fractions:def}{pecahannya}. Soal ini mengembangkan
teknik penjumlahan dan pembandingan bab ini
(\cref{met:g8:fractions:add}) menjadi sebuah cara --- yang diterbitkan
Fibonacci pada tahun $1202$ --- yang menghasilkan tulisan Mesir bagi
\omterm{def:g6:fractions:def}{pecahan} apa pun antara $0$ dan $1$, lalu berakhir pada teka-teki
warisan yang terkenal.

\medskip
\noindent\textbf{Bagian I --- \omterm{def:g6:fractions:def}{Pecahan} satuan dan kesamaan
pemecahannya.}
\begin{enumerate}
  \item Hitunglah $\frac12 + \frac13$ dan $\frac12 + \frac14$, lalu
        simpulkan tulisan Mesir bagi $\frac56$ dan $\frac34$.
  \item Periksalah catatan juru tulisnya untuk $\frac27$: tunjukkan
        bahwa $\frac14 + \frac{1}{28} = \frac27$.
  \item Periksalah bahwa $\frac13 + \frac16 = \frac12$ dan bahwa
        $\frac14 + \frac{1}{12} = \frac13$. Lalu buktikan
        \emph{kesamaan pemecahan}: untuk setiap bilangan bulat
        $n \geq 1$,
        \[
        \frac{1}{n} = \frac{1}{n+1} + \frac{1}{n(n+1)} .
        \]
  \item Simpulkan sebuah tulisan Mesir bagi $1$ itu sendiri, sebagai
        \omterm{def:g1:addition:def}{jumlah} tiga \omterm{def:g6:fractions:def}{pecahan} satuan yang berbeda-beda.
  \item Pecahlah suku terakhir $\frac56 = \frac12 + \frac13$ untuk
        memperoleh tulisan Mesir yang \emph{kedua} bagi $\frac56$,
        dengan tiga \omterm{def:g6:fractions:def}{pecahan} satuan. Jelaskan mengapa tidak ada
        \omterm{def:g6:fractions:def}{pecahan} yang tulisan Mesirnya tunggal.
\end{enumerate}

\medskip
\noindent\textbf{Bagian II --- Cara rakus Fibonacci.}
Caranya adalah: kurangkan dengan \omterm{def:g6:fractions:def}{pecahan} satuan \emph{terbesar} yang
muat, lalu ulangi pada \omterm{def:g3:division:remainder}{sisanya}.
\begin{enumerate}[resume]
  \item Dengan \omterm{def:g4:fractions:def}{penyebut} yang sama, tunjukkan bahwa
        $\frac14 < \frac27 < \frac13$. Mengapa itu membuktikan bahwa
        $\frac14$ adalah \omterm{def:g6:fractions:def}{pecahan} satuan terbesar yang lebih kecil
        daripada $\frac27$?
  \item Hitunglah $\frac27 - \frac14$. Tulisan $\frac27$ yang mana
        yang dihasilkan cara rakus itu?
  \item Sekarang ambillah $\frac57$. Tunjukkan bahwa \omterm{def:g6:fractions:def}{pecahan} satuan
        terbesar yang lebih kecil daripada $\frac57$ adalah
        $\frac12$, lalu hitunglah $\frac57 - \frac12$.
  \item Tunjukkan bahwa \omterm{def:g6:fractions:def}{pecahan} satuan terbesar yang lebih kecil
        daripada $\frac{3}{14}$ adalah $\frac15$ (bandingkan
        $\frac{3}{14}$ dengan $\frac14$ dan dengan $\frac15$),
        hitunglah $\frac{3}{14} - \frac15$, lalu simpulkan:
        \[
        \frac57 = \frac12 + \frac15 + \frac{1}{70} .
        \]
        Periksalah tulisan ini secara langsung dengan \omterm{def:g4:fractions:def}{penyebut} yang
        sama $70$.
  \item Jalankan cara rakusnya pada $\frac45$, dengan memeriksa pada
        tiap langkahnya bahwa \omterm{def:g6:fractions:def}{pecahan} satuanmu adalah yang terbesar
        yang muat, lalu periksalah tulisan akhirnya.
\end{enumerate}

\medskip
\noindent\textbf{Bagian III --- Mengapa caranya selalu berhenti, dan
unta yang kedelapan belas.}
\begin{enumerate}[resume]
  \item Pada pertanyaan 7--10 caranya diterapkan pada $\frac27$,
        $\frac57$ dan $\frac45$. Untuk masing-masing, sebutkan
        \omterm{def:g4:fractions:def}{pembilang} \omterm{def:g3:division:remainder}{sisa} yang berturut-turut (sebelum penyederhanaan
        apa pun). Apa yang kamu amati?
  \item Untuk \omterm{def:g6:fractions:def}{pecahan} $\frac ab$ dan bilangan bulat $n$, tunjukkan
        bahwa
        \[
        \frac{a}{b} - \frac{1}{n}
        = \frac{a \times n - b}{b \times n} .
        \]
  \item Andaikan cara rakusnya mengurangkan $\frac1n$ dari
        $\frac ab$: maka $\frac1n$ muat tetapi $\frac{1}{n-1}$ sudah
        terlalu besar, $\frac{1}{n-1} > \frac ab$. Dengan menaruh
        keduanya pada \omterm{def:g4:fractions:def}{penyebut} yang sama $b \times (n - 1)$,
        simpulkan bahwa $b > a \times n - a$, jadi \omterm{def:g4:fractions:def}{pembilang} yang
        baru memenuhi
        \[
        a \times n - b < a .
        \]
        Jelaskan mengapa itu membuktikan bahwa caranya selalu
        berhenti.
  \item Sebuah kisah lama: seorang ayah meninggal dengan mewariskan
        $17$ ekor unta, dengan $\frac12$ kawanannya untuk anak
        sulungnya, $\frac13$ untuk anak keduanya, $\frac19$ untuk
        anak ketiganya --- dan tidak satu pun bagiannya berupa
        bilangan bulat unta. Seorang tetangga meminjamkan seekor unta
        kedelapan belas kepada keluarganya; anak-anaknya mengambil
        $9$, $6$ dan $2$ ekor unta, lalu unta pinjamannya
        dikembalikan. Hitunglah
        $\frac12 + \frac13 + \frac19$ lalu jelaskan triknya: mengapa
        kawanannya dapat dibagi dalam unta yang utuh, dan mengapa
        tidak ada anak yang menerima kurang daripada yang dijanjikan
        wasiatnya?
  \item Perbaikilah wasiatnya: usulkan tiga \omterm{def:g6:fractions:def}{pecahan} satuan yang
        \emph{berbeda-beda} yang berjumlah tepat $1$, lalu sebutkan
        besar kawanan yang terkecil yang membuat ketiga bagiannya
        berupa bilangan bulat unta.

\end{enumerate}
\end{problem}
